\subsection{整数之间的严格不等式}

命题2. 设 \(a\) 与 \(b\) 是两个整数。则 \(a < b\) 当且仅当存在一个整数 \(c > 0\) 使得 \(b = a + c\) .

如果 \(a < b\) ，则存在一个基数 \(c \leqslant b\) （使得 \(c\) 是一个整数（ \(\S 4\) ，第2号，命题2））使得 \(b = a + c\) ( \(\S 3\) ，第6号，命题13）；若 \(a \neq b\) ，我们必须有 \(c \neq 0\) 。反之，若 \(b = a + c\) 并且 \(c \neq 0\) ，那么 \(c \geqslant 1\) ，因此 \(a < a + 1 \leqslant a + c = b\) .

命题3。设 \((a_i)_{i \in I}\) 与 \((b)_{i \in I}\) 是两个整数有限族，使得 \(a_i \leqslant b_i\) 对于每一个 \(i \in I\) 并且 \(a_i < b_i\) 对于至少一个指标 \(i\) 。然后

\[
\sum_ {i \in \mathbf {I}} a _ {i} <   \sum_ {i \in \mathbf {I}} b _ {i}.
\]

如果还有 \(b_{i} > 0\) 对于每一个 \(i \in \mathbf{I}\) ，那么

\[
\prod_ {i \in I} a _ {i} <   \prod_ {i \in I} b _ {i}.
\]

让 \(j\) 是一个指标使得 \(a_{j} < b_{j}\) ，并且让 \(J = I - \{j\}\) 。然后

\[
b _ {j} = a _ {j} + c _ {j}
\]

其中 \(c_{j} > 0\) （命题2），因此（§3，第6号，命题14）

\[
\sum_ {i \in \mathbf {I}} b _ {i} = a _ {j} + c _ {j} + \sum_ {i \in \mathbf {J}} b _ {i} \geqslant c _ {j} + a _ {j} + \sum_ {i \in \mathbf {J}} a _ {i} = c _ {j} + \sum_ {i \in \mathbf {I}} a _ {i}.
\]

由于 \(c_{j} > 0\) ，命题的第一部分由命题2得出。同样地

\[
\prod_ {i \in I} b _ {i} = (a _ {j} + c _ {j}) \prod_ {i \in J} b _ {i} = a _ {j}. \prod_ {i \in J} b _ {i} + c _ {j}. \prod_ {i \in J} b _ {i} \geqslant \prod_ {i \in I} a _ {i} + c _ {j}. \prod_ {i \in J} b _ {i}.
\]

由于 \(c_{j}\) 以及所有 \(b_{i}\) 是 \(\neq 0\) ，乘积 \(c_{j}\) . \(\prod_{i \in \mathbb{J}} b_{i}\) 是 \(\neq 0\) （§3，第4号，命题7）；因此命题的第二部分成立。

推论1. 设 \(a, a',\) 与 \(b\) 是整数，使得 \(a < a'\) 且 \(b > 0\) 。则 \(a^b < a'^b\) .

我们只需将 \(a^b\) 与 \(a'^b\) 表示为整数的有限族之积（§3，第5号，命题10）并应用命题3，注意到关系 \(a < a'\) 蕴含 \(a' > 0\) .

推论2。设 \(a, b\) 与 \(b'\) 是整数，使得 \(a > 1\) 且 \(b < b'\) ；则 \(a^b < a^{b'}\) .

因为存在一个整数 \(c > 0\) 使得 \(b' = b + c\) （命题2）；由于 \(c \geqslant 1\) ，我们有 \(a^c \geqslant a > 1\) ，由此 \(a^{b'} = a^b a^c > a^b\) .

推论3。设 \(a, b, b'\) 是整数（相应地，满足 \(a > 0\) 的整数）。则 \(a + b = a + b'\) （相应地， \(ab = ab'\) ）当且仅当 \(b = b'\) .

推论4. 若 \(a\) 与 \(b\) 是满足 \(a \leqslant b\) 的整数，则存在唯一的整数 \(c\) 使得 \(b = a + c\) .

\(c\) 的存在性由第3节第6号中的命题13得出，其唯一性则由上面的推论3得出。

整数 \(c\) （满足 \(b = a + c\) ，其中 \(a \leqslant b\) ）称为整数 \(b\) 与 \(a\) 的差，并记作 \(b - a\) 。容易验证，如果 \(a, b, a', b'\) 是满足 \(a \leqslant b\) 并且 \(a' \leqslant b'\) 的整数，那么

\[
(b - a) + (b ^ {\prime} - a ^ {\prime}) = (b + b ^ {\prime}) - (a + a ^ {\prime}).
\]

